% Preámbulo de los parciales · XeLaTeX · mismo estilo de los labs (rl_metodos/docs/curso/tex/preambulo.tex)
% Tipografía: Inter (cubre subíndices y griegas Unicode), con DejaVu Sans para los pocos símbolos que le faltan.
\documentclass[11pt,a4paper]{article}
\usepackage[margin=20mm,top=18mm,bottom=20mm]{geometry}
\usepackage{fontspec}
\setmainfont{Inter}[BoldFont={Inter SemiBold}, ItalicFont={Inter Italic}, BoldItalicFont={Inter SemiBold Italic}, Scale=0.95]
\setsansfont{Inter}
\setmonofont{DejaVu Sans Mono}[Scale=0.82]
\newfontfamily\dv{DejaVu Sans}[Scale=0.9]
\usepackage{unicode-math}
\setmathfont{latinmodern-math.otf}
\usepackage{polyglossia}\setdefaultlanguage{spanish}
\usepackage{xcolor,booktabs,longtable,array,calc,tabularx,enumitem,parskip,microtype,titlesec,fancyhdr,newunicodechar}
\usepackage[most]{tcolorbox}
\usepackage[hidelinks]{hyperref}
\emergencystretch=3em
\providecommand{\tightlist}{\setlength{\itemsep}{1pt}\setlength{\parskip}{0pt}}
\setcounter{secnumdepth}{-1}

\newunicodechar{∈}{{\dv ∈}}
\newunicodechar{∎}{{\dv ∎}}
\newunicodechar{⟨}{{\dv ⟨}}
\newunicodechar{⟩}{{\dv ⟩}}
\newunicodechar{☐}{{\dv ☐}}
\newunicodechar{✔}{{\dv ✔}}
\newunicodechar{✘}{{\dv ✘}}
\AtBeginDocument{\def\square{\mbox{\dv ☐}}}

\definecolor{tinta}{HTML}{1F2933}
\definecolor{acento}{HTML}{B7410E}
\definecolor{suave}{HTML}{F4F1EC}
\definecolor{gris}{HTML}{6B7280}
\definecolor{linea}{HTML}{D6D0C6}
\color{tinta}

\titleformat{\section}{\Large\bfseries\color{acento}}{}{0em}{}
\titleformat{\subsection}{\large\bfseries}{}{0em}{}
\titleformat{\subsubsection}{\normalsize\bfseries}{}{0em}{}
\titlespacing*{\section}{0pt}{1.4em}{0.4em}
\titlespacing*{\subsection}{0pt}{1.0em}{0.25em}
\setlist{nosep, leftmargin=1.4em, itemsep=2pt}
\renewcommand{\arraystretch}{1.3}
\setlength{\LTpre}{6pt}\setlength{\LTpost}{6pt}

% citas: el escenario inventado y las notas
\renewenvironment{quote}{\begin{tcolorbox}[colback=suave, colframe=linea, boxrule=0.5pt, arc=2pt, left=6pt, right=6pt, top=4pt, bottom=4pt]}{\end{tcolorbox}}
% código corto
\let\oldverbatim\verbatim
\let\endoldverbatim\endverbatim

% espacios para escribir
\newcommand{\hueco}[1]{\rule[-1.5pt]{\dimexpr #1\dimexpr 0.33em\relax\relax}{0.45pt}}
\newcommand{\lineasrespuesta}[1]{\par\vspace{2pt}{\color{linea}\foreach\i in {1,...,#1}{\vspace{0.75cm}\hrule height 0.5pt}}\par\vspace{4pt}}
\newcommand{\filarespuesta}{\rule{0pt}{1.25cm}}

\pagestyle{fancy}\fancyhf{}
\renewcommand{\headrulewidth}{0pt}
\fancyfoot[C]{\small\color{gris}\docpie\ · página \thepage}

\newcommand{\cabecera}[3]{%
  {\huge\bfseries\color{acento}#1\par}\vspace{4pt}
  {\large #2\par}\vspace{2pt}
  {\small\color{gris}#3\par}\vspace{6pt}
  {\color{linea}\hrule height 0.6pt}\vspace{8pt}}

\newcommand{\docpie}{Parcial 1 · Raíces C · Crédito}
\begin{document}
\cabecera{Raíces C · ¿Cuánto te cobra de verdad el banco?}{Regla falsa y Newton con datos de la Superintendencia Financiera}{Entrega y sustentación: jueves 22 de octubre de 2026 · Prof. Carlos Orrego}
Cuando pides un crédito, el banco te dice tres cosas: cuánto te presta (P), a cuántos meses (n) y cuánto pagas cada mes (A). La tasa mensual i que esconden esos tres números cumple

\[P \;=\; A\sum_{k=1}^{n}(1+i)^{-k} \;=\; A\,\frac{1-(1+i)^{-n}}{i}\]

y \textbf{no hay manera de despejar i}: para saber cuánto te cobran de verdad hay que buscar una raíz. Después la pasas a tasa efectiva anual, EA = (1 + i)¹² − 1, y la comparas con dos cosas reales: lo que cobran los bancos colombianos según la \textbf{Superintendencia Financiera} (datos abiertos en datos.gov.co) y la \textbf{tasa de usura}, el techo legal por encima del cual cobrar intereses es delito (Código Penal, art. 305).

En todo el parcial trabajamos con la función (escrita \textbf{como suma}, sin dividir por i)

\[g(i) \;=\; P - A\sum_{k=1}^{n}(1+i)^{-k}, \qquad g'(i) \;=\; A\sum_{k=1}^{n}k\,(1+i)^{-k-1}.\]

\section{Reglas del parcial}

Se entrega el \textbf{jueves 22 de octubre al empezar la clase}, escrito a mano en esta misma guía, y se sustenta ese mismo día.

\begin{itemize}
\tightlist
\item
  \textbf{Individual.} Puedes discutir ideas con compañeros, pero lo que entregas lo escribes tú, con tus palabras.
\item
  \textbf{El cuaderno de Colab} de este parcial es \texttt{parcial\_raices\_credito.ipynb}. Ábrelo desde el repositorio del curso, \textbf{github.com/cmorregof/methods\_exam}, carpeta \texttt{raices-C-credito}, con el botón «Open in Colab», y corre las celdas en orden. Antes de cada celda, escribe tu predicción en esta guía: la predicción escrita después de correr no cuenta.
\item
  \textbf{Puedes usar IA, y te audito.} Si la usas, llena la bitácora del final: qué le preguntaste, qué te respondió y qué corregiste o comprobaste tú. Usar IA no baja la nota. Lo que la baja es entregar algo que no sabes explicar.
\item
  \textbf{Sustentación (jueves, 10 minutos).} Se sortea un punto de tu parcial y lo explicas en el tablero, sin la guía en la mano, con una variante (otro dato, otro valor inicial, otro nodo). La nota escrita se multiplica por un factor: 1 si lo explicas y lo reproduces; 0,6 si lo explicas a medias; 0,2 si no puedes explicar lo que entregaste.
\item
  \textbf{Puntaje.} Parte I: 12 puntos. Parte II: 14. Parte III: 12. Parte IV: 12. Total: 50. La nota es puntos ÷ 10.
\item
  \textbf{Qué cuenta.} Una respuesta correcta sin justificación vale poco. Una justificación clara con un error de cuenta vale mucho. Si algo te parece mal planteado, dilo y explica por qué: eso también se califica.
\end{itemize}

\section{Parte I · Completar una demostración (12 puntos)}

\subsection{I.1 · La tasa existe y es única (5 puntos)}

\textbf{Teorema.} Si P \textgreater{} 0, A \textgreater{} 0, n ≥ 1 y nA \textgreater{} P (pagas en total más de lo que te prestaron), entonces g tiene \textbf{exactamente una} raíz i* en el intervalo {[}0, A/P{]}.

\textbf{Demostración (llena los espacios).}

\emph{Existencia.} g es continua en {[}0, A/P{]} porque es una suma finita de funciones \hueco{22} (½ pt).

\begin{itemize}
\tightlist
\item
  En i = 0: cada término (1 + 0)⁻ᵏ vale \hueco{6} , así que g(0) = P − \hueco{6} , que es \textless{} 0 porque \hueco{14} (1 pt).
\item
  En i = A/P: para i \textgreater{} 0 la suma es geométrica y
\end{itemize}

\[\sum_{k=1}^{n}(1+i)^{-k} \;=\; \frac{1-(1+i)^{-n}}{i} \;<\; \frac{1}{i}
\qquad \text{porque } (1+i)^{-n} \;\_\_\_\_\_\_\; 0 .\]

Con i = A/P queda A·Σ \textless{} A·(1/i) = A·(P/A) = \hueco{6} , y por tanto g(A/P) \hueco{6} 0 (1 pt).

Por el teorema del \hueco{22} existe i* en (0, A/P) con g(i*) = 0 (½ pt).

\emph{Unicidad.} g′(i) = A·Σ k(1 + i)⁻ᵏ⁻¹ es una suma de términos \hueco{14} para i \textgreater{} −1, luego g′(i) \textgreater{} 0 y g es estrictamente \hueco{14} . Si hubiera dos raíces i₁ \textless{} i₂, el teorema de \hueco{14} daría un ξ entre ellas con g′(ξ) = 0: contradicción (1½ pts). ∎

\textbf{Pregunta corta (½ pt).} ¿En qué renglón se usó la hipótesis nA \textgreater{} P? Si nA = P (crédito ``sin intereses''), ¿cuál es la raíz?

\lineasrespuesta{3}

\subsection{I.2 · Por qué la regla falsa se queda con un extremo fijo (4 puntos)}

Derivando otra vez: g″(i) = −A·Σ k(k + 1)(1 + i)⁻ᵏ⁻² \textless{} 0. O sea, g es \textbf{creciente y cóncava} (se curva hacia abajo).

\textbf{Teorema.} Sea g continua y dos veces derivable en {[}a, b{]}, con g(a) \textless{} 0 \textless{} g(b), g′ \textgreater{} 0 y g″ \textless{} 0. Sea c la raíz de la recta que pasa por (a, g(a)) y (b, g(b)) (el paso de la regla falsa). Entonces c \textgreater{} i*, y por eso la regla falsa \textbf{nunca mueve el extremo a}.

\textbf{Demostración (llena los espacios).} La recta es el polinomio de Lagrange P₁ de grado 1 con nodos a y b. Por el teorema del error de interpolación, para cada x en (a, b) existe ξ en (a, b) con

\[g(x) - P_1(x) \;=\; \frac{g''(\xi)}{2}\,(x-a)(x-b).\]

Para a \textless{} x \textless{} b: (x − a) es \hueco{6} , (x − b) es \hueco{6} y g″(ξ) es \hueco{6} (usa positivo/negativo). Entonces g(x) − P₁(x) \hueco{6} 0, es decir, la recta está \hueco{14} (encima / debajo) de la curva (1½ pts).

Evaluando en x = i*: P₁(i*) \textless{} g(i*) = \hueco{6} . Como P₁ es creciente (su pendiente (g(b) − g(a))/(b − a) es positiva) y P₁(c) = 0 \textgreater{} P₁(i*), concluimos que c \hueco{6} i* (1 pt).

Entonces g(c) \hueco{6} 0, y la regla falsa reemplaza el extremo \hueco{6} (a / b). El extremo \hueco{6} nunca cambia, y todas las rectas salen del mismo punto (a, g(a)) (1 pt). ∎

\textbf{Pregunta corta (½ pt).} ¿En qué renglón se usó g″ \textless{} 0? Si g fuera convexa (g″ \textgreater{} 0) con g(a) \textless{} 0 \textless{} g(b), ¿cuál extremo se quedaría fijo?

\lineasrespuesta{3}

\subsection{I.3 · Puente con interpolación: tres métodos, dos polinomios (3 puntos)}

\begin{enumerate}
\def\labelenumi{\alph{enumi})}
\tightlist
\item
  (1 pt) La regla falsa y la secante toman como siguiente aproximación la raíz del polinomio de \textbf{Lagrange de grado 1} por dos puntos (x₀, g(x₀)) y (x₁, g(x₁)). Escribe esa raíz:
\end{enumerate}

\[x_2 \;=\; x_1 - g(x_1)\,\frac{\_\_\_\_\_\_\_\_\_\_\_\_}{\_\_\_\_\_\_\_\_\_\_\_\_}\]

¿En qué se diferencian la regla falsa y la secante al escoger los dos puntos del paso siguiente? Si las dos arrancan con los mismos puntos x₀ = 0 y x₁ = A/P, ¿por qué su \textbf{primer} iterado es el mismo?

\begin{enumerate}
\def\labelenumi{\alph{enumi})}
\setcounter{enumi}{1}
\item
  (1 pt) Newton toma la raíz del \textbf{Taylor de grado 1} alrededor de iₖ: P(i) = g(iₖ) + g′(iₖ)(i − iₖ). En el lenguaje de la unidad de interpolación, es el polinomio osculante con un nodo x₀ = iₖ y m₀ = \hueco{6} . Escribe iₖ₊₁.
\item
  (1 pt) \textbf{A mano:} en i = 0 se tiene g(0) = P − nA y g′(0) = A·(1 + 2 + ⋯ + n) = A·n(n + 1)/2. Muestra que el primer paso de Newton desde i₀ = 0 es
\end{enumerate}

\[i_1 \;=\; \frac{2\,(nA-P)}{A\,n(n+1)}\]

y calcúlalo para P = 10.000.000, n = 36, A = 368.297,50 (el crédito de la Parte II). i₁ ≈ \hueco{14} . ¿Queda a la izquierda o a la derecha de la tasa verdadera i* ≈ 0,016119?

\lineasrespuesta{3}

\begin{center}\rule{0.5\linewidth}{0.5pt}\end{center}

\section{Parte II · Predecir → correr → explicar con datos reales (14 puntos)}

Abre \texttt{parcial\_raices\_credito.ipynb} en Colab. \textbf{La Celda 1} consulta la API de datos.gov.co (Superintendencia Financiera) y, si falla, usa una copia congelada del 15 de octubre de 2026. Algunos de los datos (corte semanal del \textbf{31 de julio de 2026}; tasa efectiva anual promedio, \textbf{ponderada por monto desembolsado} en todos los créditos de ese producto en la semana):

\begin{longtable}[]{@{}
  >{\raggedright\arraybackslash}p{(\columnwidth - 4\tabcolsep) * \real{0.3333}}
  >{\raggedright\arraybackslash}p{(\columnwidth - 4\tabcolsep) * \real{0.3333}}
  >{\raggedright\arraybackslash}p{(\columnwidth - 4\tabcolsep) * \real{0.3333}}@{}}
\toprule\noalign{}
\begin{minipage}[b]{\linewidth}\raggedright
Producto
\end{minipage} & \begin{minipage}[b]{\linewidth}\raggedright
Entidad
\end{minipage} & \begin{minipage}[b]{\linewidth}\raggedright
Tasa E.A. reportada
\end{minipage} \\
\midrule\noalign{}
\endhead
\bottomrule\noalign{}
\endlastfoot
Libre inversión & Bancolombia & 21,15 \% \\
Libre inversión & Banco de Bogotá & 23,47 \% \\
Libre inversión & Banco Unión (la más alta) & 28,76 \% \\
Vivienda no VIS en pesos & Fondo Nacional del Ahorro (la más baja) & 12,01 \% \\
Vivienda no VIS en pesos & Bancolombia & 15,76 \% \\
Tarjeta de crédito (personas y empresas), compras a 25--36 meses & Bancolombia & 28,72 \% \\
Tarjeta de crédito (personas y empresas), compras a 25--36 meses & Banco Unión (la más alta) & 28,77 \% \\
\end{longtable}

\textbf{Tasa de usura} (consumo y ordinario) = 1,5 × interés bancario corriente certificado por la Superfinanciera: \textbf{28,79 \% E.A. en julio de 2026} (Resolución 0965) y \textbf{28,59 \% E.A. en octubre de 2026} (Resolución 1472).

(Ojo: el cuaderno imprime con punto decimal. Las tasas reportadas no incluyen seguros ni comisiones.)

En cada ítem: \textbf{(a) predice en papel ANTES de correr, (b) corre la celda, (c) explica.} Si tu predicción falla no pierdes puntos por fallar: los pierdes si no explicas por qué.

\subsection{II.1 · Del banco a la cuota y de vuelta a la tasa --- Celda 3 (3 puntos)}

\begin{enumerate}
\def\labelenumi{\alph{enumi})}
\item
  \textbf{Predice:} con la tasa de Bancolombia en libre inversión (21,15 \% E.A.), ¿de cuánto será la cuota de un crédito de 10 millones a 36 meses? Ordena regla falsa, secante y Newton de menos a más \textbf{iteraciones}, y luego de menos a más \textbf{evaluaciones} (Newton gasta 2 por iteración: g y g′).
\item
  \textbf{Corre} la Celda 3 y anota: i mensual = \hueco{14} , A = \hueco{14} pesos. Iteraciones/evaluaciones: regla falsa \hueco{6} / \hueco{6} , secante \hueco{6} / \hueco{6} , Newton \hueco{6} / \hueco{6} . ¿Recuperaste la tasa publicada?
\item
  \textbf{Explica:} la celda dice que un centavo de cuota equivale a un cambio en i de ≈ 1,6·10⁻⁹. ¿La tolerancia 10⁻¹⁰ en i es más o menos exigente que ``acertar la cuota al centavo''? Cambia \texttt{tol} a \texttt{1.5e-9} (≈ un centavo): ¿cuántas evaluaciones ahorra cada método?
\end{enumerate}

\subsection{II.2 · El extremo fijo y el orden de convergencia --- Celda 4 (4 puntos)}

\begin{enumerate}
\def\labelenumi{\alph{enumi})}
\item
  \textbf{Predice:} usando I.2, ¿qué extremo de la regla falsa se mueve? ¿En qué columna esperas una constante (\texttt{razon\_lineal} = eₖ₊₁/eₖ o \texttt{razon\_cuadratica} = eₖ₊₁/eₖ²) para la regla falsa? ¿Y para Newton?
\item
  \textbf{Corre} la Celda 4 y anota: la \texttt{razon\_lineal} de la regla falsa se estabiliza en \hueco{6} ; la constante teórica C = 1 − (pendiente de la tangente en i*)/(pendiente de la cuerda desde (0, g(0))) da \hueco{6} ; la \texttt{razon\_cuadratica} de Newton se estabiliza en \hueco{6} y la constante \textbar g″(i*)\textbar/(2g′(i*)) da \hueco{6} . En la gráfica, ¿de qué punto salen todas las cuerdas?
\item
  \textbf{Explica} (2 pts): cambia a \texttt{producto\ =\ "vivienda"}, \texttt{n\ =\ 360} (Bancolombia, 30 años). ¿Cuántas iteraciones necesita ahora la regla falsa? ¿Cuánto vale C? Con la figura, explica por qué un plazo largo empeora tanto a la regla falsa: ¿qué le pasa a g(0) = P − nA cuando n = 360? Para la secante, la \texttt{razon\_cuadratica} no se estabiliza: ¿qué orden esperarías según Burden §2.3?
\end{enumerate}

\subsection{II.3 · El crédito ``sin intereses'' y la indeterminación 0/0 --- Celda 5 (4 puntos)}

La fórmula cerrada a(i) = (1 − (1 + i)⁻ⁿ)/i y la suma a(i) = Σ (1 + i)⁻ᵏ son la misma función para i ≠ 0.

\begin{enumerate}
\def\labelenumi{\alph{enumi})}
\item
  (1 pt) \textbf{Predice:} ¿cuál es el límite de (1 − (1 + i)⁻ⁿ)/i cuando i → 0? (Pista: L'Hôpital, o el binomio (1 + i)⁻ⁿ ≈ 1 − ni.) ¿Cuántas cifras correctas crees que da la fórmula cerrada en i = 10⁻⁶, 10⁻¹⁰ y 10⁻¹⁵? ¿Y en i = 10⁻¹⁶?
\item
  (1 pt) \textbf{Corre} la Celda 5 y anota el error relativo de la fórmula cerrada en i = 10⁻⁶ \hueco{6} , 10⁻¹⁰ \hueco{6} , 10⁻¹⁵ \hueco{6} , y qué da en i = 10⁻¹⁶ \hueco{6} .
\item
  (1 pt) \textbf{Explica} con lo visto en punto flotante: ¿qué operación produce cancelación? ¿Por qué en i = 10⁻¹⁶ sale exactamente 0? (Piensa en 1 + 10⁻¹⁶ con ε ≈ 2,2·10⁻¹⁶.) ¿Por qué la forma de suma no sufre nada de esto?
\item
  (1 pt) Un almacén ofrece 1.200.000 pesos en 12 cuotas de 100.000. ¿Cuál es la tasa? Con \texttt{tol0\ =\ 1e-12}, ¿qué forma converge con Newton y cuál no? Mira los iterados de la fórmula cerrada: ¿por qué se quedan ``brincando'' cerca de 10⁻¹⁰ en vez de llegar a 0? ¿Y por qué desde i₀ = 0 la fórmula cerrada ni siquiera arranca?
\end{enumerate}

\subsection{II.4 · Una oferta: ¿es legal? --- Celda 6 (3 puntos)}

Escenario hipotético: te ofrecen 10.000.000 a 36 cuotas de 400.000.

\begin{enumerate}
\def\labelenumi{\alph{enumi})}
\item
  \textbf{Predice:} ¿la tasa efectiva anual queda por encima o por debajo de la usura de octubre (28,59 \%)? ¿Por mucho o por poco?
\item
  \textbf{Corre} la Celda 6 y anota: EA de la oferta = \hueco{14} \%, diferencia con la usura \hueco{14} puntos, cuota máxima legal \hueco{14} pesos. ¿Cuántos pesos al mes se pasa la oferta?
\item
  \textbf{Explica:} con dEA/di ≈ 15,1, ¿qué error en i te permitiría equivocarte de veredicto? ¿Alcanzaría una tolerancia de 10⁻⁴? Mira también el final de la salida: las tasas más altas de los bancos reales en julio (28,76 \% y 28,77 \%) contra la usura de julio (28,79 \%). ¿Qué te dice eso de dónde se ubican los bancos?
\end{enumerate}

\begin{center}\rule{0.5\linewidth}{0.5pt}\end{center}

\section{Parte III · El agente (12 puntos)}

\begin{quote}
\textbf{Escenario inventado para este parcial.} El agente no existe; los números de la tabla sí: los calculé con experimentos reales sobre las tasas de la Superfinanciera, y los puedes auditar con la Celda 7.
\end{quote}

Un agente de aprendizaje por refuerzo recibe un crédito (P, n, A) y escoge una acción:
\textbf{RF} = regla falsa en {[}0, A/P{]}; \textbf{Sec} = secante con x₀ = 0, x₁ = A/P; \textbf{N₀} = Newton con i₀ = 0; \textbf{N\textsubscript{A}/P} = Newton con i₀ = A/P. Todas paran cuando dos iterados consecutivos difieren menos de 10⁻¹⁰.

\textbf{Recompensa} = −(número de evaluaciones de g y g′). Newton cuesta 2 por iteración; regla falsa y secante, 1 por iteración más las 2 iniciales.

Malla de entrenamiento: P = 10 millones, con la tasa real de \textbf{cada entidad} del corte del 31/07/2026 y varios plazos. Costo promedio por crédito (menor es mejor):

\begin{longtable}[]{@{}
  >{\raggedright\arraybackslash}p{(\columnwidth - 14\tabcolsep) * \real{0.1250}}
  >{\raggedright\arraybackslash}p{(\columnwidth - 14\tabcolsep) * \real{0.1250}}
  >{\raggedright\arraybackslash}p{(\columnwidth - 14\tabcolsep) * \real{0.1250}}
  >{\raggedright\arraybackslash}p{(\columnwidth - 14\tabcolsep) * \real{0.1250}}
  >{\raggedright\arraybackslash}p{(\columnwidth - 14\tabcolsep) * \real{0.1250}}
  >{\raggedright\arraybackslash}p{(\columnwidth - 14\tabcolsep) * \real{0.1250}}
  >{\raggedright\arraybackslash}p{(\columnwidth - 14\tabcolsep) * \real{0.1250}}
  >{\raggedright\arraybackslash}p{(\columnwidth - 14\tabcolsep) * \real{0.1250}}@{}}
\toprule\noalign{}
\begin{minipage}[b]{\linewidth}\raggedright
Producto (créditos)
\end{minipage} & \begin{minipage}[b]{\linewidth}\raggedright
Tasas E.A.
\end{minipage} & \begin{minipage}[b]{\linewidth}\raggedright
Plazos (meses)
\end{minipage} & \begin{minipage}[b]{\linewidth}\raggedright
RF
\end{minipage} & \begin{minipage}[b]{\linewidth}\raggedright
Sec
\end{minipage} & \begin{minipage}[b]{\linewidth}\raggedright
N₀
\end{minipage} & \begin{minipage}[b]{\linewidth}\raggedright
N\textsubscript{A}/P
\end{minipage} & \begin{minipage}[b]{\linewidth}\raggedright
\textbf{Acción aprendida}
\end{minipage} \\
\midrule\noalign{}
\endhead
\bottomrule\noalign{}
\endlastfoot
Vivienda no VIS (16 × 4 = 64) & 12,01 -- 18,00 \% & 120, 180, 240, 360 & 29,42 & \textbf{6,75} & 12,97 & 8,53 & Secante \\
Libre inversión (32 × 5 = 160) & 16,57 -- 28,76 \% & 12, 24, 36, 60, 84 & 15,02 & \textbf{7,78} & 10,57 & 10,81 & Secante \\
Tarjeta 25--36 m (27 × 4 = 108) & 18,37 -- 28,77 \% & 6, 12, 24, 36 & 12,13 & \textbf{7,97} & 9,48 & 11,50 & Secante \\
\end{longtable}

Ninguna acción falló en ninguno de los 332 créditos.

\subsection{III.1 · ¿Verdadero o falso? (6 puntos, 1 por fila)}

Para cada afirmación del agente escribe \textbf{V} o \textbf{F} y justifica en una o dos líneas con la teoría (no basta con ``lo corrí''; usa las Celdas 3 y 7 para auditar).

\begin{longtable}[]{@{}
  >{\raggedright\arraybackslash}p{(\columnwidth - 6\tabcolsep) * \real{0.0606}}
  >{\raggedright\arraybackslash}p{(\columnwidth - 6\tabcolsep) * \real{0.4242}}
  >{\raggedright\arraybackslash}p{(\columnwidth - 6\tabcolsep) * \real{0.0909}}
  >{\raggedright\arraybackslash}p{(\columnwidth - 6\tabcolsep) * \real{0.4242}}@{}}
\toprule\noalign{}
\begin{minipage}[b]{\linewidth}\raggedright
\#
\end{minipage} & \begin{minipage}[b]{\linewidth}\raggedright
Lo que reporta el agente
\end{minipage} & \begin{minipage}[b]{\linewidth}\raggedright
V / F
\end{minipage} & \begin{minipage}[b]{\linewidth}\raggedright
Justificación
\end{minipage} \\
\midrule\noalign{}
\endhead
\bottomrule\noalign{}
\endlastfoot
A & ``En vivienda, la regla falsa me costó en promedio más de 4 veces lo que me costó la secante.'' & & \\
B & ``La regla falsa sufre en vivienda porque las tasas de vivienda son bajas: a igual plazo, entre más baja la tasa, más evaluaciones necesita.'' & & \\
C & ``Newton con i₀ = 0 nunca se pasó a la derecha de la raíz en ninguno de los 332 créditos: siempre se acercó por la izquierda.'' & & \\
D & ``Newton con i₀ = A/P nunca produjo un iterado con tasa negativa.'' & & \\
E & ``Si entreno con créditos de 100 millones en vez de 10 millones, la tabla cambia, porque g queda 10 veces más grande.'' & & \\
F & ``Como paré cuando dos iterados difieren menos de 10⁻¹⁰, la regla falsa terminó con error en i menor que 10⁻¹⁰ en todos los créditos.'' & & \\
\end{longtable}

\emph{Pistas:} para B compara en la Celda 3 el Fondo Nacional del Ahorro (12,01 \%) con el Banco Unión (18,00 \%), ambos en vivienda a 360 meses. Para D corre la Celda 7 con \texttt{producto\ =\ "tarjeta"} y sus plazos. Para F mira \texttt{error\_maximo} y recuerda la constante C de II.2.

\subsection{III.2 · ¿Por qué aprendió la secante? (3 puntos)}

\begin{enumerate}
\def\labelenumi{\alph{enumi})}
\item
  (1 pt) Newton tiene orden 2 y la secante orden (1 + √5)/2 ≈ 1,618. Pero Newton paga \textbf{2 evaluaciones} por iteración. Si cada evaluación de Newton ``vale'' medio paso, ¿cuál es su orden por evaluación? (Pista: dos evaluaciones elevan el error al cuadrado, así que una evaluación lo eleva a la potencia \hueco{6} ≈ \hueco{6}.) Compara con 1,618.
\item
  (1 pt) En vivienda, N\textsubscript{A}/P (8,53) le gana a N₀ (12,97); en tarjeta pasa al revés (11,50 contra 9,48). Usa que A/P = i/(1 − (1 + i)⁻ⁿ): ¿qué tan cerca está A/P de i cuando n = 360? ¿Y cuando n = 6? (La Celda 7 imprime el rango de (A/P − i*)/i*.)
\item
  (1 pt) Usa I.2 y la concavidad de g para explicar por qué la regla falsa es la peor en los cuatro productos aunque \textbf{nunca pierde la raíz}. ¿Qué ventaja real tiene sobre la secante que la recompensa −(evaluaciones) no está midiendo?
\end{enumerate}

\subsection{III.3 · Reward hacking y verificación (3 puntos)}

\begin{enumerate}
\def\labelenumi{\alph{enumi})}
\item
  (1 pt) Un agente ``ahorrador'' quita la división y usa f(i) = P·i − A·(1 − (1 + i)⁻ⁿ). Verifica que f(i) = i·g(i) para i ≠ 0 y calcula \textbf{f(0)}. Con la oferta de II.4 y arrancando en i₀ = 0, ¿qué tasa reporta el agente y cuántas evaluaciones gasta (Celda 7b)? ¿Por qué es la acción favorita de una recompensa −(evaluaciones)? La celda muestra además que desde i₀ = 0,001, 0,005 y 0,009804 Newton también termina en 0, y desde 0,010004 en adelante encuentra la tasa verdadera. ¿Qué tiene de especial i\textsubscript{m} ≈ 0,009904 (mira f′)?
\item
  (1 pt) Otro agente para cuando \textbar g(i)\textbar{} \textless{} 100.000 pesos (``menos de una cuota de error''). La Celda 7b dice que reporta 28,27 \% E.A.: ¡legal! Pero la tasa verdadera es 28,64 \% (ilegal). Explica con g(i) ≈ g′(i*)(i − i*) y g′(i*) ≈ 1,59·10⁸ pesos por unidad de i: ¿qué error en i permite ese criterio?
\item
  (1 pt) \textbf{El verificador de una sola evaluación.} g es creciente (I.1). Si g(i\_usura) \textless{} 0, ¿de qué lado de i\_usura está la raíz? La Celda 7b da g(i\_usura) = −5.596,18 pesos. Sin resolver ninguna ecuación, ¿es legal la oferta? ¿Qué debería medir la recompensa del agente para que no le convenga ninguno de los dos atajos de (a) y (b)?
\end{enumerate}

\begin{center}\rule{0.5\linewidth}{0.5pt}\end{center}

\section{Parte IV · Consulta: nuestra idea de paper aplicada a tu tema (12 puntos)}

Haz una búsqueda profunda (NotebookLM Deep Research u otra; anótalo en tu bitácora) sobre \textbf{el cálculo de la tasa interna de retorno (TIR) y de la tasa de un crédito}: qué métodos se usan en la práctica (Newton, secante, Brent, salvaguardas con bisección), el problema de las \textbf{múltiples TIR} cuando los flujos cambian de signo más de una vez, la \textbf{regulación} de la tasa de usura en Colombia, y si existen trabajos de aprendizaje automático o RL que escojan el método o el punto inicial.

\textbf{Entrega:}

\begin{enumerate}
\def\labelenumi{\arabic{enumi}.}
\item
  \textbf{(4 pts) Síntesis de una página, escrita por ti sin IA.} ¿Qué se sabe? ¿Qué sigue abierto? ¿Qué papel juegan el punto inicial y la forma en que se escribe la ecuación (como la raíz espuria de III.3)?
\item
  \textbf{(4 pts) Tabla de 4 fuentes auditadas} (al menos una artículo de revista con revisión por pares, una documentación oficial de software o norma, y una sobre aprendizaje automático o RL para resolver ecuaciones o escoger algoritmos):
\end{enumerate}

\begin{longtable}[]{@{}
  >{\raggedright\arraybackslash}p{(\columnwidth - 10\tabcolsep) * \real{0.1667}}
  >{\raggedright\arraybackslash}p{(\columnwidth - 10\tabcolsep) * \real{0.1667}}
  >{\raggedright\arraybackslash}p{(\columnwidth - 10\tabcolsep) * \real{0.1667}}
  >{\raggedright\arraybackslash}p{(\columnwidth - 10\tabcolsep) * \real{0.1667}}
  >{\raggedright\arraybackslash}p{(\columnwidth - 10\tabcolsep) * \real{0.1667}}
  >{\raggedright\arraybackslash}p{(\columnwidth - 10\tabcolsep) * \real{0.1667}}@{}}
\toprule\noalign{}
\begin{minipage}[b]{\linewidth}\raggedright
\#
\end{minipage} & \begin{minipage}[b]{\linewidth}\raggedright
Referencia (autores, año, título, revista o sitio)
\end{minipage} & \begin{minipage}[b]{\linewidth}\raggedright
¿Existe? (enlace o DOI que abriste)
\end{minipage} & \begin{minipage}[b]{\linewidth}\raggedright
Tipo
\end{minipage} & \begin{minipage}[b]{\linewidth}\raggedright
¿Dice lo que la IA le atribuye? (cita la frase o página)
\end{minipage} & \begin{minipage}[b]{\linewidth}\raggedright
Conexión con el curso (regla falsa, secante, Newton, punto inicial, orden, punto flotante\ldots)
\end{minipage} \\
\midrule\noalign{}
\endhead
\bottomrule\noalign{}
\endlastfoot
1 \filarespuesta & & & & & \\
2 \filarespuesta & & & & & \\
3 \filarespuesta & & & & & \\
4 \filarespuesta & & & & & \\
\end{longtable}

\emph{Puntos de partida (audítalos tú, no los copies):} la documentación de la función \texttt{RATE}/\texttt{TASA} de Excel (¿qué método, qué tolerancia, cuántas iteraciones, qué punto inicial por defecto?); la regla de los signos de Descartes y la literatura sobre unicidad de la TIR (por ejemplo, la condición de Norstrøm en el \emph{Journal of Financial and Quantitative Analysis}); \texttt{scipy.optimize.brentq}.

\begin{enumerate}
\def\labelenumi{\arabic{enumi}.}
\setcounter{enumi}{2}
\tightlist
\item
  \textbf{(2 pts) Responde dos de estas tres preguntas guía:}

  \begin{itemize}
  \tightlist
  \item
    Excel pide un punto inicial (``estimar'', 10 \% por defecto) y se rinde tras unas pocas iteraciones. Con lo que viste en III.2, ¿es buen punto inicial 10 \% mensual para un crédito de vivienda a 30 años? ¿Qué propondrías tú (A/P, la fórmula de I.3c, otra)?
  \item
    Para un crédito con cuotas iguales la tasa es única (Parte I). ¿Qué cambia cuando los flujos cambian de signo varias veces (un proyecto con inversión, ganancias y luego un costo de cierre)? ¿Cuántas raíces puede haber y qué hace un método como la regla falsa o Brent en ese caso?
  \item
    ¿Existen trabajos que aprendan a escoger el método o el punto inicial (selección de algoritmos, ``learning to optimize'', RL)? ¿Lo hacen para TIR o tasas, o para problemas más generales?
  \end{itemize}
\item
  \textbf{(2 pts) Cierre.} Nuestra idea de paper: agentes de RL que aprenden las mismas decisiones que un estudiante de métodos numéricos (qué método usar, dónde empezar, cuándo parar). \textbf{¿Qué tendría que ser verdad para que nuestra idea fuera nueva en tu tema?} Sé concreto: ¿qué no ha hecho nadie según tus fuentes, y cómo lo comprobarías? ¿Qué recompensa usarías para que el agente no caiga en los atajos de III.3?
\end{enumerate}

\section{Bitácora de IA}

Una fila por cada uso. Si no usaste IA, escribe «No usé IA» y firma.

\begin{longtable}[]{@{}
  >{\raggedright\arraybackslash}p{(\columnwidth - 6\tabcolsep) * \real{0.1047}}
  >{\raggedright\arraybackslash}p{(\columnwidth - 6\tabcolsep) * \real{0.2674}}
  >{\raggedright\arraybackslash}p{(\columnwidth - 6\tabcolsep) * \real{0.3140}}
  >{\raggedright\arraybackslash}p{(\columnwidth - 6\tabcolsep) * \real{0.3140}}@{}}
\toprule\noalign{}
\begin{minipage}[b]{\linewidth}\raggedright
Punto
\end{minipage} & \begin{minipage}[b]{\linewidth}\raggedright
Qué le pedí
\end{minipage} & \begin{minipage}[b]{\linewidth}\raggedright
Qué me respondió, en una línea
\end{minipage} & \begin{minipage}[b]{\linewidth}\raggedright
Qué comprobé o corregí yo, y cómo
\end{minipage} \\
\midrule\noalign{}
\endhead
\bottomrule\noalign{}
\endlastfoot
\filarespuesta & & & \\
\filarespuesta & & & \\
\filarespuesta & & & \\
\filarespuesta & & & \\
\filarespuesta & & & \\
\filarespuesta & & & \\
\filarespuesta & & & \\
\end{longtable}

\textbf{Firma:} \hueco{30}

\section{Sustentación (la llena el profesor el jueves)}

\textbf{Punto sorteado:} \hueco{14}

\begin{itemize}
\tightlist
\item[$\square$]
  Explica el punto con sus palabras, sin leer la guía
\item[$\square$]
  Reproduce en el tablero la cuenta o el paso clave
\item[$\square$]
  Responde una variante: otro dato, otro valor inicial, otro nodo
\item[$\square$]
  Lo que dice coincide con lo que entregó y con su bitácora
\end{itemize}

\textbf{Factor:} 1 ☐ 0,6 ☐ 0,2 ☐ \textbf{Nota escrita:} \hueco{6} / 50 \textbf{Nota final:} \hueco{6}

\end{document}
